\appendix

% Los anexos y la bibliografia quedan fuera del computo de extension, asi que se numeran en
% una secuencia propia A1, A2, ... en lugar de continuar la del cuerpo: el folio del cuerpo
% termina en su ultima pagina real y la conformidad se lee sin restar nada.
%
% Por que un prefijo y no \pagenumbering{Roman}, que seria lo natural: babel-spanish
% compone los folios romanos en versalitas (\es@scroman), de modo que los anclajes de
% hyperref de los preliminares ya se expanden a page.I, page.II en mayuscula. Una secuencia
% Roman aqui chocaria frontalmente con ellos y multiplicaria los avisos "destination with
% the same identifier", que ya existen por otra causa (la portada corre en arabigo 1-2 antes
% del \pagenumbering{roman} y colisiona con el 1-2 del cuerpo). Con el prefijo, los anclajes
% son page.A1 ... page.A10 y no pueden colisionar con nada. En el PDF esto se traduce en un
% tercer rango de /PageLabels con prefijo /P, que es exactamente lo que el formato preve: el
% lector de PDF muestra "A3" en su casilla de pagina, igual que el folio impreso.
\clearpage
\setcounter{page}{1}
\renewcommand{\thepage}{A\arabic{page}}
\renewcommand{\paginastotal}{\pageref{LastPage}}

\chapter{Catálogo de variables}\label{cap:anexo1}

\begin{sloppypar}
Este anexo recoge, con fines de consulta autocontenida, la taxonomía completa de las
161 variables predictoras del modelo, obtenidas mecánicamente mediante
\texttt{build\_features.feature\_columns()} sobre la matriz procesada
\texttt{data/\allowbreak{}processed/\allowbreak{}features\_daily.parquet} y organizadas en los mismos seis grupos
temáticos empleados a lo largo de toda la memoria (en el análisis exploratorio del
capítulo \ref{cap:datos_features}, en la ingeniería de características de la sección
\ref{sec:ingenieria_features} y en el análisis de importancia de variables del capítulo
\ref{cap:resultados}), de modo que cualquier cifra de importancia de variable citada en
esta memoria pueda contrastarse directamente contra el grupo y el recuento exacto de su
familia sin necesidad de releer el código fuente.

Toda variable de este catálogo lleva el prefijo \texttt{feat\_}, una convención que no
es meramente estética: el propio código deriva el conjunto de columnas del modelo a
partir de ese prefijo, de forma que la matriz de entrada nunca se define mediante una
lista mantenida a mano que pudiera desincronizarse silenciosamente del código que
realmente construye las variables. Los seis grupos son, en el orden en que aparecen en
el cuadro \ref{tab:anexo1_grupos_features}: los siete retardos de la demanda total
(\texttt{feat\_total\_lag\_\{1,2,3,7,14,21,28\}}); los 28 retardos de operador, siete por
cada uno de los cuatro operadores del CRTM
(\texttt{feat\_\allowbreak{}\{metro,emt,carretera,cercanias\}\_\allowbreak{}lag\_\{1,2,3,7,14,21,28\}}), presentes en el modelo pese a que
\texttt{total} sea el único objetivo de predicción, en cumplimiento de la directriz del
contexto técnico del proyecto de no podar los operadores como ``columnas sin uso''; las
cuatro
medias y desviaciones típicas móviles de la demanda total, ventanas de 7 y 28 días,
todas construidas con desplazamiento previo estricto
(\texttt{feat\_total\_roll\_\{mean,std\}\_\{7,28\}}); los seis términos armónicos de
Fourier deterministas para la estacionalidad semanal de primer orden y la anual de
primer y segundo orden (\texttt{feat\_fourier\_\{weekly,annual\}\_k\{1,2\}\_\{sin,cos\}});
las once variables de calendario, que codifican de forma exhaustiva y sin nivel de
referencia implícito el tipo de día, el tipo de festividad y los indicadores de fin de
semana y de día puente; y, con diferencia el grupo más numeroso, las 105 variables
meteorológicas retardadas —21 variables físicas de Open-Meteo multiplicadas por cuatro
rezagos (1, 2, 3 y 7 días) más una media móvil de 7 días—, sujetas a la misma
prohibición de meteorología del mismo día que se documenta y verifica en el Anexo
\ref{cap:anexo2}.

\begin{table}[htbp]
  \centering
  \begin{tabular}{>{\raggedright\arraybackslash}p{4.3cm}>{\centering\arraybackslash}p{2.0cm}>{\raggedright\arraybackslash}p{7.8cm}}
\toprule
\rowcolor{naranja}
\textbf{Grupo} & \textbf{N.º variables} & \textbf{Descripción} \\
\midrule
\rowcolor{filaclara} Retardos de la demanda total & 7 & lags 1, 2, 3, 7, 14, 21, 28 \\
Retardos de operadores & 28 & 4 operadores x 7 retardos \\
\rowcolor{filaclara} Medias/\allowbreak{}desv. móviles de la demanda & 4 & media y desv. típica, ventanas 7 y 28, shift-first \\
Términos de Fourier & 6 & semanal k=1, anual k=1 y k=2 \\
\rowcolor{filaclara} Calendario & 11 & day\_\allowbreak{}type (5), holiday\_\allowbreak{}type (4), is\_\allowbreak{}weekend, is\_\allowbreak{}bridge\_\allowbreak{}day \\
Meteorología retardada & 105 & 21 variables x (lags 1/\allowbreak{}2/\allowbreak{}3/\allowbreak{}7 + media móvil 7) \\
\rowcolor{filaclara} Total & 161 &  \\
\bottomrule
\end{tabular}

  \caption{\headlinecolor{Resumen de los seis grupos de variables del modelo, 161 en total (repetición de la tabla \ref{tab:grupos_features} para consulta autocontenida del anexo).}}
  \label{tab:anexo1_grupos_features}
\end{table}

El detalle de variable individual está en los cuadros de top-15 variables por ganancia de
cada modelo XGBoost del capítulo \ref{cap:resultados}, y la lista nominal completa de las
161 columnas la genera \texttt{src/features/build\_features.py}, consultable en el
repositorio del proyecto.
\end{sloppypar}

\chapter{Suite de tests y garantías anti-fuga}\label{cap:anexo2}

\begin{sloppypar}
La afirmación central de esta memoria —que un resultado negativo bien instrumentado es
una contribución metodológica válida— solo se sostiene si el pipeline que lo produce
está libre de los errores de fuga de información que sistemáticamente producen
resultados optimistas y no reproducibles en la literatura de predicción de series
temporales. Este anexo documenta la arquitectura de verificación que respalda esa
afirmación: \textbf{470 tests}, recolectados y ejecutados con
\texttt{python -m pytest} desde la raíz del repositorio con el entorno
\texttt{.\textbackslash venv\textbackslash Scripts\textbackslash python.exe} activado,
de los cuales 257 se acumulan a lo largo de las fases de modelado (0 a 6b, cuadro
\ref{tab:cronograma_fases}) y el resto se añaden durante la redacción de la memoria:
auditan la capa de presentación (figuras y cuadros), las tres fases diagnósticas de las
secciones \ref{sec:anatomia_residuo_etapa1}, \ref{sec:paridad_informativa} y
\ref{sec:desglose_subgrupos}, el cuaderno principal de flujo, y la coherencia interna del
propio documento: \texttt{test\_prose\_claims\_match\_artifacts.py} fija contra su
artefacto de origen cada cifra de alta relevancia citada en la prosa —métricas de test,
tamaños de partición, recuentos de filas y variables— junto con la lista completa de
ficheros en que cada una debe aparecer, de modo que una cifra no pueda quedar obsoleta en
un capítulo mientras se actualiza en otro. Tres pruebas más fijan mecánicamente que las
tres fases diagnósticas, puramente explicativas, no introducen ningún modelo nuevo en la
tabla maestra de comparación ni escriben artefacto alguno bajo \texttt{models/}.

Una incidencia de la última revisión merece registro porque ilustra el modo de fallo que
esta suite existe para evitar. Las dos guardias anti-fuga del ensamblado eran
\textbf{inalcanzables}: comparaban el bloque de variables contra el nombre \emph{desnudo}
de la columna (\texttt{metro}, \texttt{temperature\_2m\_mean}), que
\texttt{feature\_columns()} descarta por construcción al filtrar por prefijo, de modo que
la lista de infractoras nunca podía tener elementos. Ninguna fuga llegó a producirse —la
protección real es estructural: los constructores solo emiten formas retardadas y
móviles—, pero una comprobación que promete una garantía que no puede dar es peor que no
tenerla, porque el siguiente lector la da por buena. Ahora comparan contra el nombre
\emph{prefijado}, que es lo que produciría una regresión de constructor, y disparan.

La suite no se limita a comprobar que el código se ejecuta sin excepciones: una parte
sustancial de ella está diseñada específicamente para hacer fallar el pipeline si se
reintrodujera, ahora o en una modificación futura, alguna de las siguientes reglas
anti-fuga, listadas de la más general a la más específica del dominio de este proyecto.

\begin{enumerate}
  \item \textbf{Desplazamiento previo estricto en todo retardo o estadístico móvil.}
    Toda media o desviación típica móvil derivada de \texttt{total} o de un operador
    aplica \texttt{shift(1)} \emph{antes} de la ventana
    (\texttt{s.shift(1).rolling(w).mean()}, nunca la corrección a posteriori ni una
    ventana sin desplazar), de forma que ninguna fila conozca su propio valor dentro de
    su propia estadística. Verificado en \texttt{tests/\allowbreak{}test\_features.py},
    con una salvedad documentada: para una ventana \emph{trailing} ambos órdenes son
    matemáticamente idénticos, y la regla se exige por corrección estructural.
  \item \textbf{Particiones siempre cronológicas, nunca aleatorias.} Ninguna partición
    se genera mediante barajado; todas derivan de la única función
    \texttt{src.utils.splits.chronological\_split}, con sus fronteras fijadas por
    \texttt{test\_split\_\allowbreak{}boundaries\_\allowbreak{}are\_\allowbreak{}exact\_\allowbreak{}and\_\allowbreak{}stable}
    en \texttt{tests/\allowbreak{}test\_splits.py}, incluida la protección contra el error
    de redondeo que un \texttt{int()} ingenuo introduciría al calcular el 70\,\% de
    1.310 filas.
  \item \textbf{Escalado ajustado exclusivamente sobre entrenamiento.} Todo
    \texttt{MinMaxScaler} de la etapa LSTM se ajusta (\texttt{fit}) solo sobre el
    entrenamiento de su propio fold y se aplica (\texttt{transform}) sin reajuste sobre
    validación y test, nunca mediante un \texttt{fit\_transform} sobre el conjunto
    completo; verificado en \texttt{tests/\allowbreak{}test\_lstm.py}.
  \item \textbf{Prohibición de ingeniería de características global antes de la
    partición}, salvo la única excepción sancionada: los términos de Fourier son una
    función determinista de la fecha, no una estadística estimada a partir del objetivo,
    de modo que calcularlos sobre la serie completa equivale a calcularlos por partición;
    fijado por
    \texttt{test\_fourier\_\allowbreak{}is\_\allowbreak{}deterministic\_\allowbreak{}and\_\allowbreak{}split\_\allowbreak{}invariant}.
  \item \textbf{Predicciones de la etapa 1 siempre fuera de muestra (\emph{out-of-fold})
    para alimentar la etapa 2.} Ningún residuo que entrena el XGBoost de la etapa 2
    procede de una predicción \emph{in-sample} de la LSTM; el esquema
    \emph{walk-forward} expansivo de 5 folds y la exclusión documentada del fold 1
    degenerado se verifican en \texttt{tests/\allowbreak{}test\_lstm.py} y
    \texttt{tests/\allowbreak{}test\_hybrid\_residual.py}. Es la decisión de diseño fijada
    por escrito como no negociable con anterioridad a cualquier resultado.
  \item \textbf{Guardia de fuga \emph{same-day} sobre la suma de operadores.} Puesto
    que \texttt{metro + emt + carretera + cercanias == total} de forma exacta, un
    operador del mismo día no es variable predictiva sino el objetivo expresado de otra
    forma; \texttt{lag\_features.\_assert\_no\_same\_day\_leakage} hace fallar la
    construcción si alguno apareciera sin retardo, comprobado sobre la matriz ensamblada
    y no solo módulo a módulo.
  \item \textbf{Guardia de fuga \emph{same-day} sobre meteorología.} Ninguna de las 21
    variables meteorológicas del día predicho entra en la matriz; solo sus versiones
    retardadas (1, 2, 3, 7 días y media móvil de 7), verificado por
    \texttt{weather\_features.\_assert\_no\_same\_day\_weather} con el mismo criterio de
    doble comprobación que la guardia de operadores.
  \item \textbf{El modelo de control \texttt{xgboost\_alone} entrena sobre las 889
    filas completas de entrenamiento, no sobre las 552 del conjunto residual.} No
    depende de los folds \emph{out-of-fold} de la LSTM, y limitar sus datos por una
    razón ajena a su diseño penalizaría artificialmente al control frente al híbrido;
    fijado por
    \texttt{test\_xgboost\_\allowbreak{}alone\_\allowbreak{}trains\_\allowbreak{}on\_\allowbreak{}the\_\allowbreak{}full\_\allowbreak{}889\_\allowbreak{}row\_\allowbreak{}split}.
  \item \textbf{El veredicto del experimento A de sensibilidad se evalúa en código, no
    a ojo.} El criterio pre-registrado de la sección \ref{sec:experimento_a} (mejora en
    días irregulares \emph{y} coste global acotado al 5\,\%) se calcula mediante
    \texttt{evaluate\_criterion\_a}, de forma que el veredicto ``NO CUMPLIDO'' no pueda
    derivar a posteriori hacia una lectura más favorable; verificado en
    \texttt{tests/\allowbreak{}test\_sensitivity\_analysis.py}.
  \item \textbf{El tema visual institucional es un objeto compartido, no una paleta
    duplicada.} El cuadro de mando interactivo y la exportación estática importan
    literalmente el mismo objeto \texttt{VIU\_TEMPLATE}, no dos definiciones que
    pudieran divergir en silencio; verificado por
    \texttt{test\_interactive\_\allowbreak{}and\_\allowbreak{}static\_\allowbreak{}paths\_\allowbreak{}share\_\allowbreak{}one\_\allowbreak{}template\_\allowbreak{}object}
    en \texttt{tests/\allowbreak{}test\_theme.py}. No es anti-fuga en sentido estricto,
    pero cumple para la capa de presentación el mismo papel de garantía estructural.
\end{enumerate}

Cada una de estas diez comprobaciones incluye, además de su verificación sobre datos
reales, al menos una prueba negativa que inyecta deliberadamente la condición de fuga
correspondiente sobre un caso de prueba mínimo y comprueba que el código
\emph{sí} falla; una guardia que nunca se ha visto obligada a disparar en la práctica es
indistinguible de una guardia que no existe, y este proyecto no se conforma con esa
ambigüedad.
\end{sloppypar}

\chapter{Reproducibilidad y entorno}\label{cap:anexo3}

\begin{sloppypar}

Este anexo documenta lo estrictamente necesario para regenerar, desde una instalación
limpia, cualquier cifra, tabla o figura de esta memoria, sin dejar ninguna decisión de
entorno implícita en la memoria de quien ejecutó el pipeline originalmente.

\textbf{Repositorio.} Todo el código de este trabajo ---ingesta, ingeniería de
características, modelos, evaluación, cuadernos y las fuentes \LaTeX{} de esta
memoria--- es de acceso público en \url{https://github.com/AdayCuestaCorrea/TFM}. Las
menciones a ``el repositorio'' en los anexos de esta memoria se refieren
siempre a esa dirección. Los datos de partida no se versionan allí: proceden de tres
fuentes abiertas cuyo origen, formato y convenciones de lectura detalla el cuadro
\ref{tab:fuentes_datos}, y hay que descargarlos y colocarlos sin renombrar en
\texttt{data/\allowbreak{}raw/} antes del primer comando de la secuencia siguiente, que
es el único punto del pipeline que los lee.

\textbf{Entorno de ejecución.} Python 3.13.5, en un entorno virtual dedicado
(\texttt{./\allowbreak{}venv}), con todas las dependencias fijadas por versión exacta, nunca por
rango, en \texttt{requirements.txt}, entre ellas:

{\small
\begin{itemize}[nosep]
  \item \texttt{numpy==2.2.6}, \texttt{pandas==2.3.3}, \texttt{scikit-learn==1.7.2}
  \item \texttt{xgboost==3.0.5}, \texttt{tensorflow==2.20.0}, \texttt{statsmodels==0.14.6}
  \item \texttt{plotly==6.3.0}, \texttt{pyarrow==22.0.0}, \texttt{kaleido==1.3.0}
\end{itemize}
}

Dos pines documentan una incompatibilidad real y no una preferencia: \texttt{pyarrow},
porque \texttt{pandas} necesita un motor Parquet explícito, y \texttt{kaleido} en su rama
1.x y nunca en la 0.2.1, porque esta cuelga de forma indefinida y sin excepción al
exportar figuras bajo \texttt{plotly} 6.x ---un fallo silencioso, más costoso de
diagnosticar que un error de instalación.

\textbf{Fijación de semillas aleatorias.} La semilla global \texttt{GLOBAL\_SEED = 42} se
define una sola vez, en \texttt{src/\allowbreak{}utils/\allowbreak{}seed.py}, y desde ahí se
aplica a \texttt{random}, \texttt{numpy} y \texttt{tensorflow} en todo punto de entrada
que entrena; ningún módulo de \texttt{src/\allowbreak{}models/} la redefine localmente.
Además de sembrar los generadores, \texttt{set\_global\_seed} exige a TensorFlow núcleos
deterministas (\texttt{TF\_\allowbreak{}DETERMINISTIC\_\allowbreak{}OPS=1} y
\texttt{tf.config.\allowbreak{}experimental.\allowbreak{}enable\_\allowbreak{}op\_\allowbreak{}determinism()}):
sembrar por sí solo no garantiza resultados bit a bit idénticos si el motor conserva
reducciones no deterministas, y es la reproducibilidad bit a bit la que este trabajo se
fijó como convención.

\textbf{Comandos de reproducción de principio a fin.} Partiendo de un clon del
repositorio, el entorno se crea con \texttt{python -m venv venv} y se puebla con
\texttt{pip install -r requirements.txt}.
Ejecutados después en este orden desde la
raíz del repositorio, con el entorno virtual activado
(\texttt{.\allowbreak{}\textbackslash{}venv\allowbreak{}\textbackslash{}Scripts%
\allowbreak{}\textbackslash{}Activate.ps1} en PowerShell), regeneran de forma
determinista la totalidad de los artefactos citados en esta memoria, sin reentrenar ni
repredecir ningún modelo más allá de lo que cada paso indica explícitamente:

{\small
\begin{enumerate}[nosep]
  \item \texttt{python -m src.ingestion.unify} — unifica las tres fuentes crudas en
    \texttt{data/\allowbreak{}interim/\allowbreak{}unified\_daily.parquet}.
  \item \texttt{python -m src.features.build\_features} — construye la matriz de 161
    variables en \texttt{data/\allowbreak{}processed/\allowbreak{}features\_daily.parquet}.
  \item \texttt{python -m src.models.baselines.run\_baselines} — baselines y SARIMAX
    seleccionado por AIC (~6 min, dominados por la rejilla de búsqueda).
  \item \texttt{python -m src.models.lstm.window\_comparison} — ventanas
    $W \in \{14, 28, 60\}$ sobre validación;
    \texttt{lstm\_window\_comparison.parquet} (tabla \ref{tab:ventana_lstm}).
  \item \texttt{python -m src.models.lstm.oof} — \emph{walk-forward} expansivo de 5 folds
    (sección \ref{sec:esquema_oof}); \texttt{lstm\_oof\_predictions.parquet} y
    \texttt{lstm\_oof\_fold\_log.parquet}.
  \item \texttt{python -m src.models.lstm.final\_model} — etapa 1 final sobre las 889
    filas; \texttt{lstm\_val\_test\_predictions.parquet} y
    \texttt{models/\allowbreak{}lstm\_stage1\_final.keras}.
  \item \texttt{python -m src.models.hybrid\_residual.xgboost\_residual} — etapa 2 sobre
    el conjunto residual de 552 filas.
  \item \texttt{python -m src.models.hybrid\_residual.xgboost\_alone} — modelo de control
    sobre las 889 filas de entrenamiento.
  \item \texttt{python -m src.models.hybrid\_residual.combine} — combina las etapas 1 y 2
    en la predicción híbrida final.
  \item \texttt{python -m src.evaluation.full\_comparison} — tabla maestra de comparación;
    \texttt{full\_comparison.parquet}, contra la que se verifican la ablación de bloques y
    el cuaderno principal de flujo.
  \item \texttt{python -m src.models.hybrid\_residual.xgboost\_residual\_weighted} —
    experimento A (sección \ref{sec:experimento_a}): peso triple en días de calendario
    irregular.
  \item \texttt{python -m src.evaluation.ensemble\_baseline} — experimento B (sección
    \ref{sec:mecanismo_ensemble}): ensamble SARIMAX + XGBoost-alone, sin reentrenamiento.
  \item \texttt{python -m src.evaluation.sensitivity\_report} — consolida ambos y evalúa
    en código el criterio pre-registrado del experimento A;
    \texttt{sensitivity\_comparison.parquet}.
  \item \texttt{python -m src.evaluation.lstm\_learning\_curve} — sección
    \ref{sec:anatomia_residuo_etapa1}: 20 ajustes de la etapa 1 (6 fracciones x 3
    semillas, más las semillas 3 y 4 en la fracción 1,00, base del control de la sección
    \ref{sec:paridad_lstm}), solo validación; \texttt{lstm\_learning\_curve.parquet}.
  \item \texttt{python -m src.evaluation.stage1\_residual\_anatomy} — tablas de la
    anatomía del residuo y veredictos de los cuatro criterios pre-registrados.
  \item \texttt{python -m src.evaluation.feature\_block\_ablation} — sección
    \ref{sec:paridad_informativa}: \texttt{xgboost\_alone} sobre cuatro subconjuntos de
    variables; \texttt{feature\_block\_ablation.parquet}.
  \item \texttt{python -m src.evaluation.subgroup\_breakdown} — sección
    \ref{sec:desglose_subgrupos}: reanálisis por quince subgrupos pre-registrados, con
    \emph{bootstrap} de bloques móviles y corrección Benjamini-Hochberg;
    \texttt{subgroup\_breakdown.parquet} y \texttt{subgroup\_rolling\_mae.parquet}.
  \item \texttt{python -m src.evaluation.lstm\_parity\_control} — sección
    \ref{sec:paridad_lstm}: LSTM en paridad informativa, 10 ajustes (2 ámbitos x 5
    semillas), solo validación; \texttt{lstm\_parity\_control.parquet}.
  \item \texttt{python -m src.evaluation.export\_figures} — exporta las 34 figuras
    estáticas de \texttt{reports/\allowbreak{}figures/}, fuera de todo kernel de Jupyter.
  \item \texttt{python -m src.evaluation.export\_latex\_tables} — vuelca las 32 tablas
    LaTeX de \texttt{docs/\allowbreak{}LaTeX/\allowbreak{}tables/}.
  \item \texttt{jupyter nbconvert -{}-to notebook -{}-execute -{}-inplace
    notebooks/\allowbreak{}07\_flujo\_completo.ipynb} — ejecuta el cuaderno principal de
    flujo (~40 s) y \textbf{contrasta cada métrica recomputada contra el valor publicado}
    en \texttt{full\_comparison.parquet}: si alguna se desvía más de $10^{-6}$ en términos
    relativos, levanta una excepción.
  \item \texttt{python -m pytest} — ejecuta la suite completa de 470 tests (descrita
    en el Anexo \ref{cap:anexo2}).
  \item \texttt{latexmk -pdf -interaction=nonstopmode -halt-on-error TFT.tex}, ejecutado
    desde \texttt{docs/\allowbreak{}LaTeX/}, compila esta memoria hasta obtener \texttt{TFT.pdf}.
\end{enumerate}
}

\textbf{Cuadro de mando interactivo.} La etapa de despliegue del ciclo KDD (sección
\ref{sec:planificacion}) se materializa en el cuaderno
\texttt{notebooks/\allowbreak{}06\_dashboard.ipynb}, que se abre con \texttt{jupyter lab}
una vez ejecutados los pasos anteriores y presenta, sobre los artefactos ya persistidos,
los cinco elementos de visualización comprometidos en la propuesta inicial del trabajo:
demanda real, demanda predicha, comparación entre modelos, error de predicción e
importancia de variables del modelo XGBoost. El cuaderno lee de
\texttt{data/\allowbreak{}processed/} y de los modelos ya entrenados de
\texttt{models/}, y no entrena ni predice nada; la exportación de las figuras
estáticas de esta memoria se realiza aparte, mediante
\texttt{export\_figures}, fuera de todo kernel de Jupyter.

\textbf{Cuaderno principal de flujo.} El cuaderno
\texttt{notebooks/\allowbreak{}07\_flujo\_completo.ipynb} recorre el trabajo completo en las
seis etapas del ciclo —carga y comprobación de datos, preparación de variables, partición
temporal, ejecución de los modelos principales, comparación de resultados y gráficos
finales—, importando de \texttt{src/} sin duplicar lógica. Se distingue del cuadro de mando
anterior en propósito y en método: aquel visualiza artefactos ya persistidos; este
\textbf{reejecuta el pipeline}, recomputando en vivo la ingesta y sus validaciones, la
matriz de 161 variables, la partición cronológica, los cuatro baselines, el SARIMAX —del
que solo se carga el orden $(2,1,2)\times(0,1,2,7)$ ya seleccionado por AIC— y los dos
modelos XGBoost. La etapa 1 LSTM es el único modelo que se carga de artefacto, por una
razón metodológica y no de coste: la etapa 2 debe corregir los residuos del mismo modelo
que produjo los resultados aquí reportados, y reentrenar la red los calcularía contra un
modelo distinto del que documenta esta memoria. Se versiona \emph{con sus salidas
guardadas}, de forma que pueda leerse sin ejecutarlo, y no escribe nada en \texttt{data/}
ni en \texttt{models/}.

\textbf{Identidad visual institucional.} El color naranja empleado en la memoria es una
\textbf{aproximación} documentada del color corporativo de la universidad, no el valor
oficial ---no disponible durante el desarrollo del proyecto---, y se codifica literalmente
en \textbf{exactamente dos puntos}, uno por capa: la capa de figuras en
\texttt{VIU\_COLORS[\textquotedbl{}primary\textquotedbl{}] = \textquotedbl{}\#E8590C\textquotedbl{}}
(\texttt{src/\allowbreak{}evaluation/\allowbreak{}theme.py}) y la capa \LaTeX{} en
\texttt{docs/\allowbreak{}LaTeX/\allowbreak{}preamble.tex} (\texttt{\textbackslash{}definecolor} de \texttt{naranja} y
\texttt{slcolor}). Sustituirlo por el valor oficial es, por diseño, un cambio de dos
líneas que se propaga a todas las figuras y cuadros sin tocar ningún otro fichero.
\end{sloppypar}

\chapter{Declaración de uso de inteligencia artificial generativa}\label{cap:anexo4}

\begin{sloppypar}

Este anexo declara el uso de herramientas de inteligencia artificial generativa en la
realización de este Trabajo Fin de Máster, delimita las tareas en las que se emplearon
y deja explícito qué decisiones permanecen, sin excepción, bajo la responsabilidad del
autor.

\section{Herramientas empleadas y alcance}\label{sec:ia_herramientas}

En la realización de este Trabajo Fin de Máster se han utilizado herramientas de
inteligencia artificial generativa, en concreto Claude (Anthropic), en su interfaz
conversacional y en su versión para desarrollo de software (Claude Code), durante el
periodo abril--septiembre de 2026 \cite{anthropic2026claude,anthropic2026claudecode}.

Su uso se ha concentrado en cuatro tareas.

\textbf{Implementación del código.} La escritura de los módulos de Python del
repositorio (ingesta, ingeniería de características, modelado, evaluación y generación
de figuras y tablas), así como de la suite de 470 pruebas automáticas que verifica su
comportamiento (Anexo \ref{cap:anexo2}), se ha realizado con asistencia de IA a partir
de especificaciones técnicas definidas por el autor. Las decisiones de arquitectura del
pipeline, el protocolo \emph{out-of-fold}, las reglas anti-fuga, la partición
cronológica y el diseño de los controles diagnósticos fueron fijadas por el autor antes
de su implementación y verificadas contra los resultados obtenidos.

\textbf{Redacción de la memoria.} El texto de este documento se ha redactado con apoyo
de IA a partir de los resultados experimentales, las decisiones metodológicas y la
estructura argumental establecidos por el autor. Los borradores generados han sido
revisados y validados por el autor, que asume la responsabilidad sobre su contenido.

\textbf{Revisión y auditoría documental.} Se ha empleado asistencia de IA para verificar
la correspondencia entre las cifras citadas en la memoria y los artefactos generados por
el código, detectar inconsistencias internas y comprobar el cumplimiento de los
requisitos formales.

\textbf{Búsqueda bibliográfica.} Al inicio del desarrollo de este TFM el autor empleó
asistencia de IA para localizar de forma expeditiva las referencias del proyecto. Cada
ficha fue revisada después de forma independiente: se comprobó su existencia, se
consultó la fuente primaria y, cuando correspondía, se descargó el documento original.
Las referencias de esta memoria no son citas inventadas ni alucinadas por el modelo;
todas están verificadas contra su fuente.

\section{Tareas no delegadas}\label{sec:ia_no_delegadas}

La definición del problema y de los objetivos, la selección de las fuentes de datos, las
decisiones metodológicas ---incluidas las que fijaron criterios de aceptación antes de
observar resultados---, la interpretación de los hallazgos y la formulación de las
conclusiones son responsabilidad del autor. En particular, la decisión de mantener y
documentar el resultado negativo de la arquitectura híbrida, en lugar de explorar
configuraciones alternativas hasta obtener un resultado favorable, es una decisión
metodológica del autor y constituye la aportación central del trabajo.

El autor asume la responsabilidad plena sobre la metodología, los datos, los resultados,
las referencias y el contenido final de esta memoria.

\section{Referencia de las herramientas}\label{sec:ia_referencias}

Las fichas de las herramientas empleadas, recogidas también en la bibliografía general,
son las siguientes:

\begin{itemize}[nosep]
  \item Anthropic. (2026). \emph{Claude} [modelo de lenguaje de gran tamaño].
    \url{https://claude.ai}.
  \item Anthropic. (2026). \emph{Claude Code} [herramienta de desarrollo de software
    asistido por IA]. \url{https://claude.com/claude-code}.
\end{itemize}

\end{sloppypar}


\chapter{Cuadros diagnósticos extendidos}\label{cap:anexo5}

Este anexo recoge los cuadros de detalle de las secciones \ref{sec:anatomia_residuo_etapa1},
\ref{sec:paridad_informativa} y \ref{sec:desglose_subgrupos}. El argumento, las cifras
titulares y los veredictos preregistrados permanecen íntegros en el cuerpo del capítulo
\ref{cap:resultados}; aquí solo se traslada el detalle tabular, al que el texto sigue
remitiendo por número de cuadro.

\begin{table}[htbp]
  \centering
  \small
  \begin{tabular}{>{\centering\arraybackslash}p{2.2cm}>{\centering\arraybackslash}p{1.8cm}>{\centering\arraybackslash}p{1.8cm}>{\centering\arraybackslash}p{2.2cm}>{\centering\arraybackslash}p{2.2cm}>{\centering\arraybackslash}p{2.2cm}}
\toprule
\rowcolor{naranja}
\textbf{Retardo (días)} & \textbf{ACF} & \textbf{PACF} & \textbf{Fuera de banda} & \textbf{Ljung-Box Q} & \textbf{p-valor (LB)} \\
\midrule
\rowcolor{filaclara} 1 & 0,088 & 0,088 & No & -- & -- \\
2 & 0,066 & 0,059 & No & -- & -- \\
\rowcolor{filaclara} 3 & 0,020 & 0,009 & No & -- & -- \\
7 & 0,162 & 0,153 & Sí & 19,48 & 0,0068 \\
\rowcolor{filaclara} 14 & 0,119 & 0,084 & Sí & 35,42 & 0,0013 \\
21 & 0,030 & 0,024 & No & -- & -- \\
\rowcolor{filaclara} 28 & 0,166 & 0,214 & Sí & 68,77 & 0,0000 \\
\bottomrule
\end{tabular}

  \caption{\headlinecolor{Autocorrelación (ACF) y autocorrelación parcial (PACF) del residuo de la etapa 1 en los retardos semanales, con las columnas de Ljung-Box (descriptivas; ver el texto). Serie primaria: modelo final sobre 393 días contiguos de validación y test.}}
  \label{tab:autocorrelacion_residuo}
\end{table}

\begin{table}[htbp]
  \centering
  \small
  \setlength{\tabcolsep}{4pt}
  \resizebox{\textwidth}{!}{\begin{tabular}{>{\raggedright\arraybackslash}p{1.9cm}>{\centering\arraybackslash}p{0.8cm}>{\centering\arraybackslash}p{2.2cm}>{\centering\arraybackslash}p{2.2cm}>{\centering\arraybackslash}p{2.2cm}>{\centering\arraybackslash}p{2.2cm}>{\centering\arraybackslash}p{1.8cm}}
\toprule
\rowcolor{naranja}
\textbf{Día} & \textbf{n} & \textbf{Residuo medio} & \textbf{Mediana} & \textbf{Desv. típica} & \textbf{|Residuo| medio} & \textbf{Media /\allowbreak{} sigma} \\
\midrule
\rowcolor{filaclara} Lunes & 56 & 142.660 & 283.700 & 663.544 & 479.203 & 0,257 \\
Martes & 56 & -42.271 & 44.308 & 563.530 & 346.908 & -0,076 \\
\rowcolor{filaclara} Miércoles & 56 & 45.764 & 154.731 & 517.848 & 313.560 & 0,082 \\
Jueves & 56 & -173.984 & -76.090 & 618.420 & 279.800 & -0,313 \\
\rowcolor{filaclara} Viernes & 56 & 125.534 & 276.219 & 690.421 & 489.159 & 0,226 \\
Sábado & 56 & -190.065 & -227.190 & 289.830 & 287.858 & -0,342 \\
\rowcolor{filaclara} Domingo & 57 & 48.428 & 106.565 & 341.577 & 226.810 & 0,087 \\
\bottomrule
\end{tabular}
}
  \caption{\headlinecolor{Perfil del residuo de la etapa 1 por día de la semana: media, mediana, desviación típica, error absoluto medio y media relativa a la $\sigma$ global del residuo. Serie primaria (393 días).}}
  \label{tab:residuo_dia_semana}
\end{table}

\begin{table}[htbp]
  \centering
  \small
  \setlength{\tabcolsep}{4pt}
  \resizebox{\textwidth}{!}{\begin{tabular}{>{\raggedright\arraybackslash}p{2.4cm}>{\raggedright\arraybackslash}p{5.2cm}>{\centering\arraybackslash}p{0.9cm}>{\centering\arraybackslash}p{2.6cm}>{\centering\arraybackslash}p{2.6cm}}
\toprule
\rowcolor{naranja}
\textbf{Dimensión} & \textbf{Grupo} & \textbf{n} & \textbf{|Residuo| medio} & \textbf{Error abs. medio (\%)} \\
\midrule
\rowcolor{filaclara} day\_\allowbreak{}type & domingo & 57 & 226.810 & 8,29 \\
day\_\allowbreak{}type & festivo & 14 & 1.818.580 & 78,98 \\
\rowcolor{filaclara} day\_\allowbreak{}type & laborable & 269 & 309.658 & 6,08 \\
day\_\allowbreak{}type & sabado & 53 & 268.777 & 7,69 \\
\rowcolor{filaclara} day\_\allowbreak{}type & laborable, bridge day & 15 & 826.986 & 21,34 \\
day\_\allowbreak{}type & laborable, ordinary & 254 & 279.107 & 5,17 \\
\rowcolor{filaclara} holiday\_\allowbreak{}type & festivo local de la ciudad de madrid & 2 & 1.960.171 & 54,10 \\
holiday\_\allowbreak{}type & festivo nacional & 11 & 1.931.389 & 89,82 \\
\rowcolor{filaclara} holiday\_\allowbreak{}type & festivo de la comunidad de madrid & 1 & 294.503 & 9,58 \\
holiday\_\allowbreak{}type & none & 379 & 291.481 & 6,63 \\
\bottomrule
\end{tabular}
}
  \caption{\headlinecolor{Magnitud del residuo de la etapa 1 por régimen de calendario (tipo de día y tipo de festivo). El residuo se concentra 6,52x en festivos y 2,96x en puentes frente a los laborables ordinarios. La fila ``laborable'' ($n=269$) es la suma de ``laborable, ordinary'' (254) y ``laborable, bridge day'' (15); no es un grupo adicional independiente.}}
  \label{tab:residuo_calendario}
\end{table}

\begin{table}[htbp]
  \centering
  \small
  \setlength{\tabcolsep}{4pt}
  \resizebox{\textwidth}{!}{\begin{tabular}{>{\raggedright\arraybackslash}p{3.6cm}>{\centering\arraybackslash}p{1.8cm}>{\centering\arraybackslash}p{1.6cm}>{\centering\arraybackslash}p{1.6cm}>{\centering\arraybackslash}p{1.7cm}>{\centering\arraybackslash}p{1.9cm}>{\centering\arraybackslash}p{2.4cm}}
\toprule
\rowcolor{naranja}
\textbf{Variable meteorológica} & \textbf{Retardo} & \textbf{Rho bruto} & \textbf{Rho parcial} & \textbf{p parcial} & \textbf{p ajustado (BH)} & \textbf{Significativa} \\
\midrule
\rowcolor{filaclara} snowfall\_\allowbreak{}sum & lag\_\allowbreak{}3 & -0,109 & -0,132 & 0,0090 & 0,8149 & No \\
snowfall\_\allowbreak{}sum & lag\_\allowbreak{}1 & -0,096 & -0,118 & 0,0196 & 0,8149 & No \\
\rowcolor{filaclara} snowfall\_\allowbreak{}sum & lag\_\allowbreak{}7 & 0,112 & 0,114 & 0,0244 & 0,8149 & No \\
relative\_\allowbreak{}humidity\_\allowbreak{}2m\_\allowbreak{}max & lag\_\allowbreak{}1 & 0,283 & 0,106 & 0,0372 & 0,8149 & No \\
\rowcolor{filaclara} weather\_\allowbreak{}code & lag\_\allowbreak{}7 & 0,189 & 0,100 & 0,0489 & 0,8149 & No \\
wind\_\allowbreak{}speed\_\allowbreak{}10m\_\allowbreak{}min & lag\_\allowbreak{}7 & 0,080 & 0,097 & 0,0567 & 0,8149 & No \\
\rowcolor{filaclara} wind\_\allowbreak{}speed\_\allowbreak{}10m\_\allowbreak{}max & lag\_\allowbreak{}2 & -0,147 & -0,094 & 0,0630 & 0,8149 & No \\
wind\_\allowbreak{}speed\_\allowbreak{}10m\_\allowbreak{}mean & lag\_\allowbreak{}7 & 0,030 & 0,092 & 0,0699 & 0,8149 & No \\
\rowcolor{filaclara} rain\_\allowbreak{}sum & lag\_\allowbreak{}1 & 0,169 & 0,090 & 0,0774 & 0,8149 & No \\
temperature\_\allowbreak{}2m\_\allowbreak{}max & lag\_\allowbreak{}2 & -0,280 & -0,090 & 0,0776 & 0,8149 & No \\
\bottomrule
\end{tabular}
}
  \caption{\headlinecolor{Diez variables meteorológicas retardadas con mayor correlación parcial (en valor absoluto) con el residuo de la etapa 1. ``Rho parcial'' controla por los cuatro términos anuales de Fourier; ninguna es significativa tras Benjamini-Hochberg (q=0,10).}}
  \label{tab:residuo_meteo}
\end{table}

\begin{table}[htbp]
  \centering
  \small
  \setlength{\tabcolsep}{4pt}
  \resizebox{\textwidth}{!}{\begin{tabular}{>{\centering\arraybackslash}p{1.4cm}>{\raggedright\arraybackslash}p{2.4cm}>{\centering\arraybackslash}p{1.6cm}>{\centering\arraybackslash}p{1.6cm}>{\centering\arraybackslash}p{1.3cm}>{\centering\arraybackslash}p{1.3cm}>{\centering\arraybackslash}p{1.6cm}>{\centering\arraybackslash}p{1.5cm}>{\centering\arraybackslash}p{1.7cm}}
\toprule
\rowcolor{naranja}
\textbf{Partición} & \textbf{Modelo} & \textbf{MAE} & \textbf{RMSE} & \textbf{MAPE (\%)} & \textbf{R²} & \textbf{Delta MAE} & \textbf{Delta MAE (\%)} & \textbf{Cuota del hueco (\%)} \\
\midrule
\rowcolor{filaclara} test & lstm\_\allowbreak{}alone & 280.952 & 431.714 & 6,62 & 0,8948 & 0 & 0,0 & 0,0 \\
test & hybrid & 234.223 & 328.131 & 5,49 & 0,9392 & -46.729 & -16,6 & 37,4 \\
\rowcolor{filaclara} test & xgboost\_\allowbreak{}alone & 156.121 & 234.154 & 3,30 & 0,9690 & -78.103 & -33,3 & 62,6 \\
val & lstm\_\allowbreak{}alone & 411.142 & 655.479 & 11,81 & 0,8023 & 0 & 0,0 & 0,0 \\
\rowcolor{filaclara} val & hybrid & 280.738 & 415.562 & 7,09 & 0,9205 & -130.405 & -31,7 & 78,4 \\
val & xgboost\_\allowbreak{}alone & 244.843 & 355.027 & 5,68 & 0,9420 & -35.895 & -12,8 & 21,6 \\
\bottomrule
\end{tabular}
}
  \caption{\headlinecolor{Cadena de error \texttt{lstm\_alone} $\rightarrow$ \texttt{hybrid} $\rightarrow$ \texttt{xgboost\_alone} por partición. ``Delta MAE'' es el paso respecto a la fila anterior (negativo = reduce el MAE); ``Cuota del hueco'' expresa ese paso como fracción de la distancia total \texttt{lstm\_alone}-\texttt{xgboost\_alone}. No es una descomposición.}}
  \label{tab:cadena_error}
\end{table}

\begin{table}[htbp]
  \centering
  \small
  \setlength{\tabcolsep}{4pt}
  \resizebox{\textwidth}{!}{\begin{tabular}{>{\centering\arraybackslash}p{1.6cm}>{\centering\arraybackslash}p{1.6cm}>{\centering\arraybackslash}p{1.7cm}>{\centering\arraybackslash}p{2.4cm}>{\centering\arraybackslash}p{2.1cm}>{\centering\arraybackslash}p{2.1cm}>{\centering\arraybackslash}p{2.2cm}>{\centering\arraybackslash}p{2.0cm}}
\toprule
\rowcolor{naranja}
\textbf{Fracción} & \textbf{Filas train} & \textbf{Secuencias} & \textbf{MAE val (mediana)} & \textbf{MAE val (mín)} & \textbf{MAE val (máx)} & \textbf{R² val (mediana)} & \textbf{Fits degenerados} \\
\midrule
\rowcolor{filaclara} 0.25 & 222 & 194 & 1.167.187 & 1.167.111 & 1.181.916 & 0,1990 & 3 \\
0.40 & 356 & 328 & 515.603 & 492.361 & 517.640 & 0,7580 & 0 \\
\rowcolor{filaclara} 0.55 & 489 & 461 & 502.449 & 492.995 & 512.022 & 0,7565 & 0 \\
0.70 & 622 & 594 & 490.773 & 449.494 & 495.200 & 0,7762 & 0 \\
\rowcolor{filaclara} 0.85 & 756 & 728 & 411.457 & 401.565 & 430.084 & 0,8033 & 0 \\
1.00 & 889 & 861 & 416.324 & 379.074 & 433.829 & 0,7985 & 0 \\
\bottomrule
\end{tabular}
}
  \caption{\headlinecolor{Curva de aprendizaje de la LSTM de etapa 1: MAE de validación mediano (mín-máx sobre 3 semillas; 5 en la fracción 1,00, que es la base de comparación del control de paridad de la sección \ref{sec:paridad_lstm}) por fracción del bloque de entrenamiento. ``Fits degenerados'' cuenta las semillas cuyo ajuste colapsó a una predicción cuasi-constante. La referencia \texttt{xgboost\_alone} (val) es 244.843.}}
  \label{tab:curva_aprendizaje}
\end{table}

\begin{table}[htbp]
  \centering
  \small
  \setlength{\tabcolsep}{4pt}
  \resizebox{\textwidth}{!}{\begin{tabular}{>{\raggedright\arraybackslash}p{2.3cm}>{\centering\arraybackslash}p{1.4cm}>{\centering\arraybackslash}p{1.0cm}>{\centering\arraybackslash}p{1.0cm}>{\centering\arraybackslash}p{1.0cm}>{\centering\arraybackslash}p{1.8cm}>{\centering\arraybackslash}p{1.8cm}>{\centering\arraybackslash}p{2.0cm}}
\toprule
\rowcolor{naranja}
\textbf{Subconjunto} & \textbf{N variables} & \textbf{prof.} & \textbf{n\_\allowbreak{}est.} & \textbf{lr} & \textbf{MAE val} & \textbf{MAE test} & \textbf{Delta MAE val (\%)} \\
\midrule
\rowcolor{filaclara} completo & 161 & 5 & 100 & 0,05 & 244.843 & 156.121 & 0,0 \\
solo temporal & 39 & 4 & 100 & 0,05 & 461.854 & 303.927 & 88,6 \\
\rowcolor{filaclara} solo exogeno & 122 & 4 & 300 & 0,05 & 320.520 & 311.023 & 30,9 \\
sin meteo & 56 & 5 & 100 & 0,05 & 237.040 & 155.196 & -3,2 \\
\rowcolor{filaclara} solo calendario & 17 & 5 & 300 & 0,10 & 300.734 & 279.830 & 22,8 \\
\bottomrule
\end{tabular}
}
  \caption{\headlinecolor{Ablación de bloques de variables sobre \texttt{xgboost\_alone} (control diagnóstico). MAE de validación y test por subconjunto, con los hiperparámetros seleccionados en cada búsqueda y el cambio relativo de MAE de validación frente a \texttt{completo}. Estudio de retirada de bloques, NO una descomposición: los bloques no son disjuntos en información y las columnas no suman a nada.}}
  \label{tab:ablacion_bloques}
\end{table}

\begin{table}[htbp]
  \centering
  \small
  \setlength{\tabcolsep}{4pt}
  \resizebox{\textwidth}{!}{\begin{tabular}{>{\raggedright\arraybackslash}p{1.8cm}>{\centering\arraybackslash}p{0.8cm}>{\centering\arraybackslash}p{2.0cm}>{\centering\arraybackslash}p{1.8cm}>{\centering\arraybackslash}p{1.8cm}>{\centering\arraybackslash}p{1.9cm}>{\centering\arraybackslash}p{1.7cm}>{\centering\arraybackslash}p{2.0cm}}
\toprule
\rowcolor{naranja}
\textbf{Ámbito} & \textbf{k} & \textbf{MAE val (mediana)} & \textbf{MAE val (mín)} & \textbf{MAE val (máx)} & \textbf{Hueco cerrado (\%)} & \textbf{Fits degenerados} & \textbf{best epoch (mediana)} \\
\midrule
\rowcolor{filaclara} calendario & 17 & 388.625 & 382.695 & 395.877 & 16,2 & 0 & 75,0 \\
completo & 161 & 363.133 & 354.235 & 391.630 & 31,0 & 0 & 53,0 \\
\bottomrule
\end{tabular}
}
  \caption{\headlinecolor{LSTM en paridad informativa por ámbito de variables exógenas (control diagnóstico, 5 semillas, solo validación). MAE de validación mediano con banda mín-máx, porcentaje del hueco univariante$\to$\texttt{xgboost\_alone} cerrado, número de ajustes degenerados y \texttt{best\_epoch} mediano. Base univariante: mediana de 5 semillas 416.324; mejor semilla 411.142.}}
  \label{tab:paridad_lstm}
\end{table}

\begin{table}[htbp]
  \centering
  \small
  \setlength{\tabcolsep}{4pt}
  \resizebox{\textwidth}{!}{\begin{tabular}{>{\raggedright\arraybackslash}p{2.7cm}>{\centering\arraybackslash}p{0.7cm}>{\raggedright\arraybackslash}p{2.3cm}>{\centering\arraybackslash}p{1.8cm}>{\raggedright\arraybackslash}p{2.1cm}>{\centering\arraybackslash}p{1.8cm}>{\centering\arraybackslash}p{1.8cm}>{\raggedright\arraybackslash}p{2.9cm}}
\toprule
\rowcolor{naranja}
\textbf{Subgrupo} & \textbf{n} & \textbf{Mejor competidor} & \textbf{MAE competidor} & \textbf{Mejor híbrido} & \textbf{MAE híbrido} & \textbf{Delta MAE} & \textbf{Veredicto} \\
\midrule
\rowcolor{filaclara} Laborable ordinario & 133 & ensemble\_\allowbreak{}equal & 138.628 & hybrid\_\allowbreak{}weighted & 200.326 & 61.699 & no gana \\
Laborable puente & 3 & persistence & 131.251 & hybrid & 399.637 & 268.386 & no inferencial (n insuficiente) \\
\rowcolor{filaclara} Sábado & 27 & xgboost\_\allowbreak{}alone & 129.381 & hybrid\_\allowbreak{}weighted & 185.918 & 56.537 & no gana \\
Domingo & 29 & xgboost\_\allowbreak{}alone & 100.200 & hybrid & 233.778 & 133.578 & no gana \\
\rowcolor{filaclara} Festivo & 5 & ensemble\_\allowbreak{}inverse\_\allowbreak{}mae & 224.587 & hybrid & 666.572 & 441.986 & no inferencial (n insuficiente) \\
Fin de semana & 56 & xgboost\_\allowbreak{}alone & 114.269 & hybrid\_\allowbreak{}weighted & 215.528 & 101.259 & no gana \\
\rowcolor{filaclara} Entre semana & 141 & ensemble\_\allowbreak{}equal & 143.426 & hybrid\_\allowbreak{}weighted & 222.176 & 78.750 & no gana \\
Calendario irregular & 8 & ensemble\_\allowbreak{}inverse\_\allowbreak{}mae & 221.011 & hybrid & 566.472 & 345.461 & no inferencial (n insuficiente) \\
\rowcolor{filaclara} Calendario ordinario & 189 & ensemble\_\allowbreak{}equal & 136.193 & hybrid\_\allowbreak{}weighted & 204.831 & 68.638 & no gana \\
Verano (jul-ago) & 33 & xgboost\_\allowbreak{}alone & 80.274 & hybrid\_\allowbreak{}weighted & 133.293 & 53.019 & no gana \\
\rowcolor{filaclara} Resto del año & 164 & ensemble\_\allowbreak{}equal & 150.371 & hybrid\_\allowbreak{}weighted & 237.791 & 87.420 & no gana \\
Periodo Semana Santa & 10 & xgboost\_\allowbreak{}alone & 332.267 & hybrid & 605.909 & 273.642 & no analizable (sin solape val-test) \\
\rowcolor{filaclara} Demanda anómala & 16 & xgboost\_\allowbreak{}alone & 302.679 & hybrid\_\allowbreak{}weighted & 520.736 & 218.057 & no inferencial (n insuficiente) \\
Demanda regular & 181 & ensemble\_\allowbreak{}equal & 119.662 & hybrid\_\allowbreak{}weighted & 193.727 & 74.065 & no gana \\
\bottomrule
\end{tabular}
}
  \caption{\headlinecolor{Desglose por subgrupo, partición test (ámbito de veredicto). ``Delta MAE'' es MAE del mejor híbrido menos MAE del mejor competidor; positivo significa que el híbrido es peor. El veredicto es único por subgrupo y se lee sobre esta partición.}}
  \label{tab:subgrupos_test}
\end{table}

\begin{table}[htbp]
  \centering
  \small
  \setlength{\tabcolsep}{4pt}
  \resizebox{\textwidth}{!}{\begin{tabular}{>{\raggedright\arraybackslash}p{2.7cm}>{\centering\arraybackslash}p{0.7cm}>{\raggedright\arraybackslash}p{2.3cm}>{\centering\arraybackslash}p{1.8cm}>{\raggedright\arraybackslash}p{2.1cm}>{\centering\arraybackslash}p{1.8cm}>{\centering\arraybackslash}p{1.8cm}>{\raggedright\arraybackslash}p{2.9cm}}
\toprule
\rowcolor{naranja}
\textbf{Subgrupo} & \textbf{n} & \textbf{Mejor competidor} & \textbf{MAE competidor} & \textbf{Mejor híbrido} & \textbf{MAE híbrido} & \textbf{Delta MAE} & \textbf{Veredicto} \\
\midrule
\rowcolor{filaclara} Laborable ordinario & 121 & ensemble\_\allowbreak{}equal & 205.597 & hybrid\_\allowbreak{}weighted & 240.264 & 34.667 & no gana \\
Laborable puente & 12 & ensemble\_\allowbreak{}inverse\_\allowbreak{}mae & 486.552 & hybrid\_\allowbreak{}weighted & 690.105 & 203.553 & no inferencial (n insuficiente) \\
\rowcolor{filaclara} Sábado & 26 & xgboost\_\allowbreak{}alone & 132.334 & hybrid\_\allowbreak{}weighted & 245.120 & 112.786 & no gana \\
Domingo & 28 & xgboost\_\allowbreak{}alone & 128.658 & hybrid & 245.450 & 116.792 & no gana \\
\rowcolor{filaclara} Festivo & 9 & ensemble\_\allowbreak{}inverse\_\allowbreak{}mae & 322.149 & hybrid & 306.409 & -15.740 & no inferencial (n insuficiente) \\
Fin de semana & 54 & xgboost\_\allowbreak{}alone & 130.428 & hybrid\_\allowbreak{}weighted & 249.490 & 119.062 & no gana \\
\rowcolor{filaclara} Entre semana & 142 & ensemble\_\allowbreak{}equal & 236.953 & hybrid\_\allowbreak{}weighted & 286.561 & 49.608 & no gana \\
Calendario irregular & 21 & ensemble\_\allowbreak{}inverse\_\allowbreak{}mae & 416.094 & hybrid & 547.929 & 131.835 & no inferencial (n insuficiente) \\
\rowcolor{filaclara} Calendario ordinario & 175 & ensemble\_\allowbreak{}inverse\_\allowbreak{}mae & 203.657 & hybrid\_\allowbreak{}weighted & 243.111 & 39.454 & no gana \\
Verano (jul-ago) & 57 & xgboost\_\allowbreak{}alone & 130.465 & hybrid\_\allowbreak{}weighted & 174.647 & 44.183 & no gana \\
\rowcolor{filaclara} Resto del año & 139 & ensemble\_\allowbreak{}inverse\_\allowbreak{}mae & 262.686 & hybrid\_\allowbreak{}weighted & 318.052 & 55.366 & no gana \\
Periodo Navidad & 16 & xgboost\_\allowbreak{}alone & 394.273 & hybrid\_\allowbreak{}weighted & 603.473 & 209.200 & no analizable (sin solape val-test) \\
\rowcolor{filaclara} Demanda anómala & 34 & xgboost\_\allowbreak{}alone & 343.987 & hybrid & 400.652 & 56.665 & no inferencial (n insuficiente) \\
Demanda regular & 162 & ensemble\_\allowbreak{}inverse\_\allowbreak{}mae & 190.627 & hybrid\_\allowbreak{}weighted & 241.797 & 51.170 & no gana \\
\bottomrule
\end{tabular}
}
  \caption{\headlinecolor{Desglose por subgrupo, partición validación (requisito de réplica). La columna ``Veredicto'' repite el veredicto preregistrado único por subgrupo; esta tabla solo difiere de la \ref{tab:subgrupos_test} en las columnas de MAE y diferencia.}}
  \label{tab:subgrupos_val}
\end{table}

\begin{table}[htbp]
  \centering
  \small
  \setlength{\tabcolsep}{4pt}
  \resizebox{\textwidth}{!}{\begin{tabular}{>{\raggedright\arraybackslash}p{2.9cm}>{\centering\arraybackslash}p{0.7cm}>{\centering\arraybackslash}p{1.9cm}>{\centering\arraybackslash}p{1.9cm}>{\centering\arraybackslash}p{1.9cm}>{\centering\arraybackslash}p{1.3cm}>{\centering\arraybackslash}p{1.3cm}>{\raggedright\arraybackslash}p{2.9cm}}
\toprule
\rowcolor{naranja}
\textbf{Subgrupo} & \textbf{n} & \textbf{Delta MAE} & \textbf{IC 95\% inf.} & \textbf{IC 95\% sup.} & \textbf{p} & \textbf{p (BH)} & \textbf{Veredicto} \\
\midrule
\rowcolor{filaclara} Laborable ordinario & 133 & 61.699 & 22.714 & 97.147 & 0,005 & 0,006 & no gana \\
Sábado & 27 & 56.537 & 8.811 & 107.801 & 0,021 & 0,021 & no gana \\
\rowcolor{filaclara} Domingo & 29 & 133.578 & 37.602 & 169.095 & 0,003 & 0,005 & no gana \\
Fin de semana & 56 & 101.259 & 31.456 & 130.511 & 0,004 & 0,005 & no gana \\
\rowcolor{filaclara} Entre semana & 141 & 78.750 & 35.267 & 120.969 & 0,000 & 0,000 & no gana \\
Calendario ordinario & 189 & 68.638 & 30.543 & 91.683 & 0,000 & 0,000 & no gana \\
\rowcolor{filaclara} Verano (jul-ago) & 33 & 53.019 & -- & -- & 0,000 & 0,000 & no gana \\
Resto del año & 164 & 87.420 & 40.550 & 116.936 & 0,000 & 0,000 & no gana \\
\rowcolor{filaclara} Demanda regular & 181 & 74.065 & 39.705 & 93.746 & 0,000 & 0,000 & no gana \\
\bottomrule
\end{tabular}
}
  \caption{\headlinecolor{Diferencia de MAE (mejor híbrido $-$ mejor competidor) con intervalo de confianza \emph{bootstrap} del 95\,\%, p bruto y p ajustado por Benjamini-Hochberg ($q=0{,}10$), para la familia inferencial en la partición test. En verano el IC se suprime porque la fracción de réplicas utilizables cae a $0{,}81$; la p se conserva.}}
  \label{tab:subgrupos_inferencia}
\end{table}
